\frontsection{Приложение А. Материалы эксперимента}
Каждому запуску назначьте отдельную папку. Ниже приведён пример для работы №\,2:
\begin{lstlisting}[language={}]
runs/B0_01/
    manifest.json
    history.csv
    validation_predictions.csv
    predictions_raw.csv
    model.keras
    preprocessing.joblib
    diagnostics.pdf
\end{lstlisting}

\file{manifest.json} фиксирует конфигурацию, метрики валидации,
идентификаторы обучающих и валидационных объектов, версии библиотек,
контрольные суммы файлов и имена подготовленных признаков.
\file{history.csv} содержит значения по эпохам.
\file{validation\_predictions.csv} связывает прогноз с известной
ценой для анализа ошибок. Модель и параметры предобработки сохраняются
вместе: применение модели к иначе подготовленным признакам некорректно.

Графики добавляются после сохранения опыта:
\begin{lstlisting}
fig.savefig("runs/B0_01/diagnostics.pdf", bbox_inches="tight")
\end{lstlisting}

Для названий файлов придерживайтесь одного идентификатора опыта:
в приведённой строке папка \file{runs/B0\_01} должна совпасть с
аргументом вашего вызова \file{save\_experiment}.
Не открывайте файлы сериализованной предобработки из неизвестных источников;
этот комплект предназначен для сохранённых вами опытов.

\subsection*{Минимальный шаблон вывода}
«При фиксированных \ldots{} изменение \ldots{} привело к \ldots{}
на валидационной части. Это видно по \ldots{}. Возможное объяснение
состоит в \ldots{}. Вывод ограничен \ldots{}. Для проверки устойчивости
потребовалось бы \ldots{}».

Заполните пропуски конкретными условиями, результатами и наблюдениями.
Фразы «модель хорошо обучилась» или «стало точнее» без указания метрики
и сравнения не объясняют полученный результат.

\subsection*{Что сохраняется в других работах}
В работе №\,1 сохраняются таблица сравнения валидации, тестовые прогнозы, classification report, обученный Pipeline и описание группового разбиения. В работе №\,3 --- список изображений с метками и разбиением, история каждой эпохи, выбранные веса, описание архитектуры, словарь классов, метрики и прогнозы.

Для повторения собственной CNN требуется её исходный код. Один файл весов не описывает архитектуру, подготовку изображения или смысл номера класса. Материалы всех трёх работ сопровождаются версиями среды, seed и графиками.

\subsection*{Словарь основных обозначений}
\begin{table}[H]\centering\small
\begin{tabularx}{\textwidth}{@{}l Y@{}}
\toprule
Символ & Значение \\
\midrule
\(x_i, y_i,\widehat y_i\) & Признаки, истинный ответ и прогноз для объекта \(i\). \\
\(d, n\) & Число признаков и число объектов в текущем наборе. \\
\(\theta\) & Обучаемые параметры модели. \\
\(B\) & Размер мини-выборки; фактический последний пакет может быть меньше. \\
\(\eta\) & Скорость обучения. \\
\(C_{in}, C_{out}\) & Число входных и выходных каналов свёртки. \\
\(H,W\) & Высота и ширина карты признаков. \\
\(TP,FP,FN\) & Верно найденные объекты класса, ложные срабатывания и пропуски. \\
\bottomrule
\end{tabularx}\end{table}
В работе №\,1 параметр \(C\) без индекса обозначает гиперпараметр LinearSVC; в матрице \(C_{ij}\) индексы обозначают пару истинного и предсказанного классов. Смысл определяется контекстом, а не одним символом.
